\section{Axiom 与 Neo Lab：AI for Math 创业需要什么}

视频描述中提到，Axiom 完成 2 亿美元 A 轮融资，估值 16 亿美元；Ken Ono 离开或暂停传统教授路径加入 Axiom 也引发关注。无论具体媒体叙事如何，这说明 AI for Math 已经从学术问题变成高资本密度 Neo Lab 问题。

\begin{figure}[H]
\centering
\includegraphics[width=0.92\textwidth]{ai-math-neolab.png}
\caption{AI for Math Neo Lab 的组成：数学家、形式化系统、模型训练、资金算力和产品目标。}
\end{figure}

\begin{knowledgebox}{读图：Neo Lab 不是普通创业公司}
图中五个要素共同指向 Axiom 这类组织：顶级数学家网络提供问题和审美，Lean/Coq 提供形式化基础，模型训练提供自动化能力，资金算力支撑规模实验，产品目标把研究转成可用系统。
\end{knowledgebox}

\subsection{为什么数学家网络重要}

数学数据不是普通网页数据。真正有价值的问题、证明路线、定理库组织和审美判断，往往掌握在数学家社区中。Ken Ono 这样的数学家加入，不只是背书，更可能带来问题选择、研究网络和验证标准。

\begin{importantbox}{AI for Math 的数据飞轮}
如果系统能帮助数学家形式化证明、补全 lemma、发现反例，数学家又能反馈修正和评价，这就可能形成数据飞轮。没有数学家参与，模型很容易停留在题库和表面证明。
\end{importantbox}

\subsection{最不可能创业的创业者}

访谈里洪乐潼谈到“最不可能创业的创业者”。这类 Neo Lab 和传统应用创业不同：创始人要理解数学、AI、形式化系统、研究组织和融资。它更像把研究院、模型团队和产品公司压缩成一个高密度组织。

\begin{warningbox}{研究导向创业的双重风险}
如果太学术，产品和收入路径不清；如果太产品化，可能失去基础研究深度。AI for Math 需要在科学目标和工程落地之间保持张力。
\end{warningbox}

\subsection{本章小结}

Axiom 代表一种新型 AI Lab：研究问题足够基础，资本密度足够高，团队需要同时连接数学家、形式化系统和模型训练。它不是普通 SaaS 创业，也不是传统学术实验室。

